\chapter{Dise\~no e implementaci\'on del sistema}

\section{Visi\'on general de la arquitectura}

El sistema se ha disenado como un workflow modular que transforma una consulta financiera estructurada en una respuesta analitica basada en calculos ejecutados sobre datos historicos. La arquitectura separa de forma explicita las fases de validacion, seleccion de agente, planificacion, generacion de codigo, ejecucion e interpretacion. Esta separacion permite razonar sobre cada etapa de forma independiente y facilita la incorporacion progresiva de modelos de lenguaje sin perder trazabilidad.

\section{Entrada estructurada y estado compartido}

La entrada del sistema no es una consulta libre sin procesar, sino una representacion normalizada que incluye la consulta original, la intencion analitica, los tickers implicados, el periodo temporal, el intervalo y las rutas a los ficheros CSV. A partir de esta entrada se construye un estado compartido que viaja por todos los nodos del workflow e incorpora campos como el agente seleccionado, el plan analitico, el codigo generado, las salidas de ejecucion, los avisos y la respuesta final.

\section{Workflow del sistema}

El flujo principal comienza con un nodo de ingesta que valida la entrada minima y comprueba la existencia de los ficheros de datos. A continuacion, un nodo de enrutado selecciona la familia analitica adecuada. El nodo especialista genera un plan estructurado, el nodo de generacion construye el script Python correspondiente, el nodo de ejecucion lanza dicho script en un proceso controlado y el nodo de interpretacion transforma la salida numerica en una explicacion legible.

\section{Agentes anal\'iticos especializados}

La logica especifica de cada intencion se organiza mediante familias de agentes. Esta organizacion evita concentrar toda la logica en un unico modulo y permite asociar cada intencion con un plan, un cuerpo de codigo y una forma de interpretar resultados. En la version actual se contemplan agentes para crecimiento y comparacion de activos, vision descriptiva, retornos, riesgo historico y analisis tecnico.

\section{Generaci\'on controlada de c\'odigo}

La generacion de codigo se plantea como una fase intermedia entre el plan analitico y la ejecucion. En el modo determinista, el sistema utiliza plantillas programaticas que producen scripts Python ajustados al contrato de salida esperado. Esta decision reduce la variabilidad durante las primeras iteraciones del MVP y permite validar primero la arquitectura antes de delegar la generacion de codigo en un modelo de lenguaje.

\section{Ejecuci\'on y artefactos generados}

El codigo generado se guarda como un script independiente y se ejecuta mediante un proceso externo controlado. La ejecucion captura la salida estandar, la salida de error, el codigo de retorno y un payload estructurado con los resultados. Estos artefactos permiten revisar que calculos se han realizado, que entrada se ha utilizado y que respuesta se ha construido a partir de la salida del script.

\section{Gesti\'on de errores y fallback determinista}

El sistema incorpora validaciones para intenciones no soportadas, rutas CSV inexistentes y restricciones de cardinalidad en los tickers. Cuando se activa la capa LLM, el diseno conserva un fallback determinista: si la llamada al modelo falla o no cumple el contrato esperado, el sistema puede continuar con la logica programatica. Esta estrategia evita que una dependencia externa comprometa la ejecucion completa del workflow.

\section{Consideraciones \'{e}ticas, seguridad y uso responsable}

El sistema se orienta exclusivamente al analisis financiero historico y descriptivo. Sus respuestas deben interpretarse como explicaciones basadas en datos pasados, no como predicciones ni como recomendaciones de inversion. Por este motivo, el diseno evita formular instrucciones de compra o venta y prioriza expresiones que expliquen metricas, periodos y limitaciones del analisis realizado.

Una primera consideracion etica es el riesgo de que el usuario confunda una interpretacion descriptiva con asesoramiento financiero. Para reducir este riesgo, la memoria y el sistema delimitan el alcance del prototipo: se trabaja con series historicas y se calculan metricas como crecimiento, retornos, volatilidad, drawdown o indicadores tecnicos basicos, pero no se realizan predicciones de rentabilidad futura ni recomendaciones personalizadas.

Una segunda consideracion se relaciona con la generacion automatica de codigo. Ejecutar codigo producido por un sistema automatico puede introducir riesgos si no existen controles adecuados. En este proyecto, la ejecucion se realiza mediante scripts acotados, con entradas estructuradas, captura de salidas y gestion explicita de errores. Ademas, la version base mantiene plantillas deterministas, lo que permite validar el flujo antes de habilitar fases generativas de mayor riesgo.

Tambien es relevante la trazabilidad. En aplicaciones financieras asistidas por modelos de lenguaje, una respuesta final aparentemente coherente puede ocultar errores de calculo o supuestos no verificados. La arquitectura propuesta mitiga este problema separando el plan, el codigo, la ejecucion y la interpretacion. De este modo, cada resultado puede relacionarse con los datos de entrada y con los artefactos generados durante el workflow.

Finalmente, el uso de datos de mercado exige prudencia. Los CSV utilizados en el proyecto tienen finalidad academica y experimental. El sistema no pretende sustituir herramientas profesionales de analisis financiero ni validar decisiones de inversion reales. Su valor principal esta en estudiar una arquitectura reproducible, modular y controlada para automatizar analisis historicos, manteniendo supervision humana sobre la interpretacion final y sobre cualquier posible evolucion del prototipo.
